\begin{claim}{F.9}{Proved}{An explicit common kernel element for the seven-point maps}
Every one of the twelve seven-point maps in F.7--F.8 has
a nontrivial kernel, both on the specified discrete
$F_{13}$ and on its full profinite completion.  With the
source coordinate, basepoint, and meridians specified
below, the same element
\[
 w=m_{-1}m_i m_1^{-1}m_i^{-1}                         \tag{F.9.1}
\]
lies in all twelve kernels.  A source winding character
sends $w$ to $1$, so its procyclic closure is a full
$\widehat\Z$ contained in their common completed kernel.
The joint kernel of all thirty-six six- and seven-point
maps in (F.7.9) is exactly this same kernel.
This refutes injectivity of these particular restriction
maps; the residual $A_5$ centralizer comparison in F.7
remains unresolved.
\end{claim}
\begin{proof}
\textbf{The specified source and target.}
Use the faithful tetrahedral action generated by
\[
 z\longmapsto-z,\qquad z\longmapsto\frac{z+i}{z-i}.
\]
Its fourteen ramification points are infinity and the
zeros of
\[
 D(z)=z(z^4-1)(z^8+14z^4+1).
\]
Thus $X=\mathbf P^1\setminus R_{A_4}$ has
$\pi_1(X,b)=F_{13}$ for $b=5+4i$.
The seven labels $I=\{0,1,2,3,4,7,10\}$ of F.8,
in that order, have the following coordinate functions:
\[
 \begin{split}
 (m_0(z),\ldots,m_6(z))=
 \left(z,\frac{i(z+1)}{z-1},-\frac1z,
 -\frac{z-i}{z+i},\frac{i(z-1)}{z+1},-z,
 \frac{z-i}{z+i}\right).
 \end{split}                                                   \tag{F.9.2}
\]
Here the subscripts $0,\ldots,6$ denote positions in
the displayed seven-set, not its original labels.
The exact finite part of the certificate below verifies
all $144$ matrix products in this $A_4$, its correspondence
with the labels in F.8, and the ramification polynomial
$D$ by multiplying the fixed-point polynomials of its
seven nontrivial cyclic subgroups.

For $z\in X$, let
\[
 L_z(t)=\frac{(t-m_0(z))(m_1(z)-m_6(z))}
              {(t-m_6(z))(m_1(z)-m_0(z))}.
\]
The seven orbit points are distinct, so this is the
unique M\"obius transformation carrying the first,
second, and seventh points to $0,1,\infty$.
Define the ordered configuration map
\[
 \Phi:X\longrightarrow\operatorname{Conf}_6(\mathbf C),
 \qquad z\longmapsto(R_0(z),\ldots,R_5(z)),
 \quad R_j(z)=L_z(m_j(z)).                              \tag{F.9.3}
\]
The rational functions extend across any removable
singularities in their displayed expressions; their
only poles are among $R_{A_4}$.  In particular
$R_0=0$ and $R_1=1$.
Appending the seventh point $\infty$ and taking the
labelled M\"obius quotient gives a map
\[
 q:\operatorname{Conf}_6(\mathbf C)\longrightarrow M_{0,I}
 \quad\hbox{with}\quad q\circ\Phi=f_I.                 \tag{F.9.4}
\]
These are actual continuous maps, with the displayed
base configurations.  The induced map of fundamental
groups therefore factors through the pure braid group
$P_6\le B_6$.  The inclusion here is induced by the
cover from ordered to unordered configurations and is
injective by path lifting.  We shall prove
$\Phi_*(w)=1$ in $B_6$, which proves $f_{I,*}(w)=1$
by this factorization.

\textbf{The three based source meridians.}
Put $r=27/1000$.  For $c=-1,i,1$, set
\[
 t_c=(\operatorname{Re}c-r)+4i,
 \quad b_c=c-r-ri,
 \quad d_c=c+r-ri,
 \quad e_c=c+r+ri,
 \quad a_c=c-r+ri.
\]
Let $m_c$ be the loop obtained by traversing the
straight polygonal segments with successive vertices
\[
 b,\ t_c,\ b_c,\ d_c,\ e_c,\ a_c,\ b_c,\ t_c,\ b.
                                                               \tag{F.9.5}
\]
The square is counterclockwise; its stem is followed
by its reverse.  Each square contains the single
ramification point $c$.  This also follows directly
from the above polynomial: the other eight finite
points have real and imaginary coordinates with
absolute values $(\sqrt3-1)/2$ or $(\sqrt3+1)/2$,
whereas the remaining finite points are
$0,\pm1,\pm i$.
The exact segment check below also verifies that every
stem and every square edge avoids all zeros of $D$.
Source loop products in (F.9.1) are chronological
concatenation.

\textbf{Exact extraction of the three braids.}
Multiply every coordinate in (F.9.3) by
$\eta=1+(173/1000)i$ and project to its real part.
This multiplication is an orientation-preserving
homeomorphism of the plane.  At the basepoint the
six coordinate positions, in increasing projected
order, are
\[
                      (0,5,3,1,4,2).
\]
Identify this chamber with the standard real ordered
base configuration by a path within the chamber.
The chamber of any fixed strict real ordering is
convex, so the identification is unique up to
homotopy.  Read crossings chronologically.  Write
$j$ for $\sigma_j$, and $-j$ for $\sigma_j^{-1}$;
the positive sign means that the strand on the left
immediately before the crossing has larger imaginary
coordinate.  This fixes both the numbering and the
crossing convention.

The executed certificate
\path{code/part_a/seven_point_kernel_f9.py}
uses only rational and Gaussian-rational arithmetic.
For a segment $z=p+(q-p)t$, $0\le t\le1$, and each
pair of normalized coordinates, write their projected
complex difference before taking real parts as
$A(t)/B(t)$ with $A,B\in\mathbf Q(i)[t]$.
It verifies that $A$, $B$, every coordinate denominator,
and $D(p+(q-p)t)$ have no complex zero on the real
segment: the real and imaginary parts have a greatest
common divisor with no real root there.
The possible crossings are precisely the roots in
$(0,1)$ of
\[
       P(t)=\operatorname{Re}(A(t)\overline{B(t)}).
                                                               \tag{F.9.6}
\]
Clearing rational denominators preserves its roots.
Sturm isolation over $\mathbf Q$ enumerates all of
them.  Their rational isolating intervals are pairwise
disjoint, their multiplicities are one, and
$\operatorname{Im}(A\overline B)$ has a certified
constant nonzero sign on each interval.  No endpoint
is a crossing.  The two strands are checked to be
adjacent in the current order at every crossing,
and the final order is checked against direct exact
evaluation at the endpoint.

There are $18$ different segments up to reversal,
$270$ pair checks, and $64$ isolated crossings in
these segments; the largest crossing-polynomial
degree is $12$.  Reversed segments give the inverse
words.  Between successive crossings the path lies
in a convex ordering chamber.  At a crossing only
the named adjacent pair changes order, with its
certified nonzero imaginary separation.  Straightening
these intervals and crossing neighborhoods gives the
usual elementary braid pieces.  Thus the recorded
words represent the actual monodromies, not just a
linear representation of them.  This is also the
projection-isotopy argument of Artin's Theorem~3
\cite{A-Artin1947}.

The freely reduced words for $m_{-1},m_i,m_1$ are,
respectively, $A,B,C$ below; a line break only separates
consecutive parts of the same word:
\[
\begin{aligned}
 A={}&(1,4,5,2,1,3,2,3,2,1,2,-5,-4,-5,\\
     &\hspace{9mm}-5,4,5,-2,-1,-2,-3,-2,-3,-1,-2,-5,-4,-1),\\
 B={}&(1,4,5,2,3,1,4,2,3,4,-3,1,2,3,4,5,\\
     &\hspace{9mm}-2,5,4,3,2,1,-4,-3,-2,-4,-1,-3,-2,-5,-4,-1),\\
 C={}&(1,4,5,2,3,4,1,2,3,4,5,-1,3,2,-5,-4,\\
     &\hspace{9mm}-4,5,-2,-3,1,-5,-4,-3,-2,-1,-4,-3,-2,-5,-4,-1).
\end{aligned}                                                   \tag{F.9.7}
\]

\textbf{The braid identity and the source detector.}
Use the faithful Artin substitution on
$F_6=\langle X_1,\ldots,X_6\rangle$:
\[
 \sigma_j(X_j)=X_jX_{j+1}X_j^{-1},\qquad
 \sigma_j(X_{j+1})=X_j,
 \qquad \sigma_j(X_k)=X_k\ (k\ne j,j+1).              \tag{F.9.8}
\]
Its inverse is used for a negative letter.  Products
of substitutions are composed from the displayed
word order; the certificate implements composition
by substitution and cancellation of adjacent inverse
letters.  Artin's Theorem~14, with the inverse
generator formula on p.~113, proves faithfulness
\cite{A-Artin1947}.  Simultaneously reversing all
generator signs is an automorphism of the braid
group, and reversing the stacking convention is an
anti-automorphism; neither changes the identity test.

For the exact words (F.9.7), substitution gives
\[
       \operatorname{Artin}(AB)=\operatorname{Artin}(BC),
 \qquad
       \operatorname{Artin}(ABC^{-1}B^{-1})(X_j)=X_j
       \quad(1\le j\le6).                              \tag{F.9.9}
\]
This is an integer free-word calculation.  The output
records all six freely reduced common images for the
first equality and all six single-letter images for
the second.  Faithfulness gives
$AB=BC$ in $B_6$.  By the certified identification of
the three words with the three source monodromies,
(F.9.1) lies in $\ker\Phi_*$, and hence in
$\ker f_{I,*}$ by (F.9.4).

The map $X\to\mathbf C^*$, $z\mapsto z+1$, gives
the integral winding character
\[
 \ell:\pi_1(X,b)\longrightarrow\Z,
 \qquad
 \ell(m_{-1})=1,\quad \ell(m_i)=\ell(m_1)=0.
\]
Therefore $\ell(w)=1$ and $w\ne1$.
Completing this exact character gives
$\widehat\ell:\widehat F_{13}\to\widehat\Z$ with
$\widehat\ell(w)=1$; this proves nontriviality in the
full completion without a detection assumption.
The continuous power map $\widehat\Z\to\widehat F_{13}$,
$a\mapsto w^a$, has composite $a\mapsto a$ with
$\widehat\ell$, so it is injective.  Its image lies
in the closed completed kernel.

Finally F.8 gives the actual map identities
$f_{hI}=r_h\circ f_I$ for the twelve left translates
$hI$, with $r_h$ a labelled target isomorphism.
They carry this same source word, and all its
profinite powers, into the identity in every target.

The twenty-four six-point conditions in (F.7.9)
follow from these twelve conditions by specified
forgetting maps.  For one of the six subgroups
$C_5\le A_5$, let $p_+,p_-$ be its fixed poles and
let $O$ be either of its two five-cycles.  Put
\[
 K=O\cup\{p_+,p_-\},\qquad J_\pm=O\cup\{p_\pm\}.
\]
Forgetting the other pole gives the labelled morphism
$d_\pm:M_{0,K}\to M_{0,J_\pm}$ and the exact identity
$f_{J_\pm}=d_\pm\circ f_K$ on the same based source.
The twelve choices of $(C_5,O)$ give all twelve
seven-sets, and their two pole choices give all
twenty-four six-sets.  The separate executed certificate
\path{code/part_a/seven_point_six_factorization_f9.py}
lists all twenty-four containments, the omitted pole,
and the six positions retained in each seven-tuple.
Functoriality of fundamental groups and of full
profinite completion proves on both sources that
\[
 \bigcap_{I\text{ among all }36}\ker f_{I,*}
       =\bigcap_{K\text{ among the }12\text{ seven-sets}}
                                      \ker f_{K,*}.     \tag{F.9.10}
\]
For the completed equality the starred maps in
(F.9.10) mean their completed homomorphisms.
Consequently (F.9.1) and its full procyclic closure
also satisfy every six-point condition.

The common-kernel injectivity route in F.7 therefore
fails.  The necessary seven-point restrictions of
an element of the residual $A_5$ centralizer still
vanish, but the converse implication has not been
proved.  No centralizer element is constructed here.

The discovery computation used numerical crossings
and a finite Burau image to find the short candidate.
Neither computation is an input to this proof.
The file
\path{code/part_a/seven_point_kernel_f9_output.json}
contains the deterministic exact certificate:
the coordinate functions, Gaussian-rational path
vertices, every crossing polynomial, every rational
root interval and sign, the three words, and the
Artin identities.  The companion stdout file records
the successful execution.
\end{proof}
